\chapter{研报写作、复核与数字护栏}
\chapterlead{模型可以写得很好，但它写下的每个数字都要有人负责。在这个系统里，负责的是护栏和数字锁。}

\section{主笔怎么写}

主笔会话拿到的材料是：真值表与财务分析、估值产物、资料包、分析框架、四五份专题笔记，以及写作规范。规范要求的结构是“结论先行”：
首页是投资要点、评级与核心数据，然后依次是年报点评、行业与竞争、深度专题、预期差、盈利预测与估值、风险提示、跟踪指标与结论失效条件。
每一节内部按“事实 → 推论 → 观点”展开，观点句写成“我们认为”“我们预计”，与事实句分开。估值部分必须写出假设、口径简化和适用范围，
全文末尾有“不构成投资建议”的声明。

写作规范里有一条我们反复强调：不要复述表格。“营业收入 1,688.38 亿元，同比下降 1.21\%；归母净利润 823.20 亿元，同比下降 4.53\%”这样的句子，
读者看图表就够了；正文应该解释这些数字为什么是这样、对明年意味着什么。茅台研报首页的写法是“茅台酒量价稳住，系列酒拖累利润”，
随后用分部毛利率和销量的变化说明这个判断。

\section{数字护栏}

文稿写完后，主笔必须运行 \code{guard}，通过后才算完成第 8 阶段。护栏做五件事。

\paragraph{出处。}正文里每个带单位的数字，都要在本次运行的结论性产物里找到对应值，或者登记在外部来源表中。
所谓“结论性产物”，是指真值表、财务分析、业绩分析、DCF 结果、可比公司结果、反向 DCF 和盈利预测摘要；三表模型的中间网格、敏感性矩阵不算，
研报断言的原文更不算，否则研报里的错误数字就会变成自己的出处。

\paragraph{按展示精度比较。}“58.62 亿元”只与真值除以 $10^8$ 后四舍五入到两位小数比较，“14.49\%”只与百分数值比较，容差是展示精度的半个末位（$\halfulp$）。
这样一来，产物里中间结果再多，也不会“碰巧”放行一个编造的数字。

\paragraph{指标与期间绑定。}只核对“有出处”还不够：把 2025 年的营业收入写在 2024 年后面，数字本身有出处，但句子是错的。
护栏识别句子里的指标名和期间，要求数字与该指标、该期间的真值一致。茅台初稿里“2025 年营业收入……同比增长 30.45\%”就是在这一步被拦下的，
因为真值表里该指标该期间的同比是 $-1.21\%$。

\paragraph{外部数字带编号。}来自来源登记表的数字，所在句子必须带对应的 \code{[S编号]}；“龙头”“第一”“份额”这类判断性措辞没有编号也会被拦。

\paragraph{前后矛盾与出处分类。}护栏调用研报内部一致性检查，比较全文各处的评级、各情景每股价值、预测年度的收入与利润是否一致。
最后，它把每个有出处的数字归类：报告披露和外部来源是\textbf{事实}，程序计算是\textbf{推论}，估值假设和模型预测是\textbf{观点}。
同一个数在多处都能找到出处时，先认事实，再认推论，最后认观点。

表~\ref{tab:wr-prov} 是 1005 批次三份终稿用当前代码重跑护栏的结果。三份研报各有 838～894 个带出处的数字，无出处、指标期间不符和外部数字缺编号均为 0；
事实类占 75\%～76\%，观点类占 15\%～18\%。

\begin{table}[htbp]
\centering
\caption{1005 批次终稿的数字出处构成}
\label{tab:wr-prov}
\small
\begin{tabular}{@{}lrrrrrrr@{}}
\toprule
\hd{公司} & \hd{事实} & \hd{推论} & \hd{观点} & \hd{合计} & \hd{无出处} & \hd{不符} & \hd{矛盾} \\
\midrule
贵州茅台 & 632 & 70 & 136 & 838 & 0 & 0 & 1（误报） \\
\rowband
宁德时代 & 648 & 80 & 126 & 854 & 0 & 0 & 0 \\
招商轮船 & 676 & 54 & 164 & 894 & 0 & 0 & 0 \\
\bottomrule
\end{tabular}
\end{table}

茅台那一条“矛盾”是护栏自己的误报，也是一个值得记录的教训。原文是“……平均目标价 1755.48 元……25 家给出‘买入’评级 [S10]”，
引述的是外部机构的评级分布，不是本报告的评级。一致性检查本来会排除带来源编号的句子，但护栏在调用它之前先删掉了全部 \code{[S编号]}，
排除规则失去了依据。修复后重跑为 0。

\begin{figure}[htbp]
\centering
\begin{tikzpicture}
\begin{axis}[xbar stacked, bar width=11pt, width=0.86\linewidth, height=4.0cm, xmin=0, xmax=920,
  symbolic y coords={招商轮船,宁德时代,贵州茅台}, ytick=data, xlabel={带出处的数字（个）},
  legend style={at={(0.5,1.04)},anchor=south}, legend columns=3, axis y line*=left, axis x line*=bottom,
  y tick label style={font=\small\heiti}]
\addplot[fill=ink!80,draw=ink] coordinates {(632,贵州茅台) (648,宁德时代) (676,招商轮船)};
\addplot[fill=gold!80,draw=gold] coordinates {(70,贵州茅台) (80,宁德时代) (54,招商轮船)};
\addplot[fill=seal!70,draw=seal] coordinates {(136,贵州茅台) (126,宁德时代) (164,招商轮船)};
\legend{事实,推论,观点}
\end{axis}
\end{tikzpicture}
\caption{三份研报中数字的出处构成。观点类数字（预测与估值假设）在研报里一律以“我们预计”“我们假设”引出。}
\label{fig:wr-prov}
\end{figure}

\subsection{护栏的误报也要修}

护栏误报不会让错误的东西流出去，但会给写作施加压力：模型倾向于删掉或模糊化那些正确但难以绑定的数字，长期看会损害报告的信息量。
2026 年 9 月 29 日茅台第二次运行护栏时剩下 3 条“不符”，其中两条是误报：护栏把“茅台酒销售公司净利率 33.35\%”绑定到公司整体净利率（50.53\%）上核对，
把“第二季度收入 387.88 亿元”绑定到全年口径上核对。当时主笔会话的做法是改写句子绕过去，功能上可行，但等于让文稿迁就护栏。
之后我们让护栏识别分部、子公司的别名和单季度的表述，按对应口径核对；1005 批次终稿里这两句以原样措辞通过。

\section{复核编辑与数字锁}

主笔写完的稿子，交给一个没有参与写作的复核编辑会话。它的技能说明把自己定位为“研究所质控编辑”，任务有三项：
删去自证、圆场和程序口吻（如“经程序校验”“根据计算结果”）；理顺故事线和论证链，检查结论是否先行、推导是否完整；
前后对账，找出同一件事在不同章节说法不一致的地方。拿不准但重要的事项，它不自己判断，而是写进人工复核清单 \file{review.md}。

复核编辑只能改措辞和结构，不能碰数字。这由数字锁 \code{numlock} 保证：改写稿里出现的数字和来源编号，必须是初稿的子集，可以删，不能新增，也不能改。
茅台研报的复核日志记录了这次改写的结果：“新增数字 0，新增来源 0，减少 23 处”。

\begin{story}[复核编辑删掉了什么：茅台研报的复核日志节选]
复核日志先回答三个问题：核心结论读完能否说出来、推导链是否完整、故事线是否连贯；然后逐节记录改动。以下是几条原样摘录的改动理由：

\begin{itemize}
  \item 小标题改为“第四季度利润是全年低点”（收入并非全年最低）；删“全年唯一一个利润大幅下滑的季度”（无依据）。
  \item 删“没有依靠大额促销”和“利润质量差距由此拉开”（五粮液上半年利润 +89.3\%，与该句矛盾）。
  \item 删“新增 1,800 吨基酒产能在 2026 年释放”（基酒需多年贮存，与护城河一节矛盾），关键假设一节同步删除。
  \item 删“悲观与乐观情景同权重”（估值只输出基准情景评级，无此权重）。
\end{itemize}

这些都是数字护栏发现不了的问题：每个数字都有出处，但论断本身站不住，或者与另一章矛盾。
\end{story}

人工复核清单是交给人看的。招商轮船的清单有 17 条，其中第 12 条是这样写的：来源 S42 称两家船东合计约 150 艘、“市场份额 40\%”，
同一来源又称其中一家约占全球 880 艘的 14\%～17\%，150/880 与 40\% 不能同时成立；正文已删去 40\%，只写 150 艘，请核对 40\% 的分母口径。
这类问题需要领域知识去判断，系统能做的是把它找出来、说清楚为什么重要、记下当前怎么处理的。

\section{排版}

排版设计师会话用排版工具 reportkit 编写一个排版脚本，把文稿、产物和外部数据组装成卖方体例的 PDF：首页是投资要点与核心数据表，
有目录、图表目录、外部资料来源附录和免责声明；图表的数据全部从产物和来源登记表读取，图下注明来源。闸门要求 25～32 页、不少于 60 个图表。
1005 批次三份交付 PDF 分别是 42、42、30 页（含目录、图表目录、外部资料来源与免责声明），83、102、70 个图表。
这里有一处我们整理材料时才发现的问题：闸门按排版预览图的张数计页，茅台和宁德时代的预览图是 25 张，闸门判通过，但交付 PDF 实际是 42 页，超出了上限。
页数闸门应当直接读 PDF，这一点列入待修事项。

排版是五个角色里第二贵的（1005 批次平均 16.9 分钟、5.83 美元），因为排版设计师要写四五百行脚本，反复渲染、逐页查看、调整版面。
10 月 3 日以前由复核会话兼做排版；v2 版把排版独立成一个 Opus 会话，按论点组织图表，中际旭创研报在盲评中“图表服务论点”一项随之由不通过变为通过。
